\subsection{Do predictive video representations preserve readable spatial structure?}

Encoder probes alone establish that a representation contains spatial information, but they do not ask whether that information survives prediction of a future latent.  We therefore distinguish three questions under one frozen ScanNet++ development protocol (612 train and 69 development scenes; horizons $h\in\{2,4\}$).  D1 trains a spatial readout on real target-view features and tests whether its accuracy is target-specific rather than a layout prior, wrong target, or copied context.  D2 trains an external future-latent predictor and compares the true camera-motion condition with identity motion, matched-horizon shuffled motion, and an independently trained context-only predictor.  D3 asks the stricter transfer question for V-JEPA2: a readout trained only on real future encoder states is frozen and applied unchanged to the model's native predicted future state.

The D1 results rule out a purely prior-based explanation.  For Wan2.1, spatial target readout gives scene-macro normalized-depth MAE $0.155\pm0.001$ at $h{=}2$ and $0.152\pm0.001$ at $h{=}4$, compared with $0.241$ and $0.238$ for the train-layout prior; paired target-minus-prior differences are $-0.086$ with 95\% scene-cluster CIs excluding zero at both horizons.  Target features likewise outperform within-scene wrong targets and copied context.  In D2, true motion improves Wan prediction over identity, shuffled, and independently trained context-only conditions at both horizons (true-minus-control MAE differences $-0.003$ to $-0.007$, all paired CIs excluding zero), showing that the pose-derived condition contributes beyond temporal context alone.

The strongest test is V-JEPA2 D3.  A spatial readout applied to real future states has MAE $0.182/0.183$ ($h{=}2/h{=}4$); without retraining, the same readout obtains $0.246/0.243$ on native predicted states, still better than copying the last context ($0.289/0.297$).  Retraining an equal-capacity reader on predicted states improves to $0.223/0.227$.  Thus prediction preserves a substantial, spatially useful signal, but changes its interface relative to the real encoder state.  This supports neither a claim of lossless spatial propagation nor one of complete spatial loss.  All results are development-only: the dataset audit found no defensible historically unused test split, and pose conditioning is not an intervention on physical action.
